% CS301 Week 10 -- File organization and hashing: sequential, direct,
% indexed-sequential and hashed files; hash-function design (division,
% multiplication, mid-square, folding, digit analysis, polynomial strings,
% universal hashing); open hashing (separate chaining) and closed hashing
% (linear/quadratic probing, double hashing); load factor and the expected
% cost of a lookup.
% Difficulty: intro (course-wide default per knowledge-base-CS301.md).
% Page-grounded in Sedgewick2011 Sec.3.4 (Hash Tables, p.458), Wirth2004
% Sec.5.1-5.4 (Key Transformations, pp.177-183), Karumanchi2017 Ch.14
% (Hashing, pp.610-639 PDF). Horowitz & Sahni Ch.8 (file organization) and
% Aho-Hopcroft-Ullman Ch.4 are cited at chapter level only (not indexed) --
% no page numbers are asserted for them.
% Handoff: Week 9 closed on "the table is sorted and searchable in O(log n);
% Week 10 asks whether we can find a record in one probe, without sorting at
% all." This deck pays that off with the calculator's 2,000-identifier table:
% O(n) as a Week 6 list, O(log n) as a Week 8 sorted array, O(1) expected as
% a hash table -- and names the price (order is lost).
% Sets up Week 11 (trees): the ordered operations hashing cannot provide.

\documentclass[aspectratio=169]{beamer}
\input{header}
\coursecode{CS301}

\title{One Probe}
\subtitle{Week 10 --- File Organization \& Hashing}
\author[Dr. Auqib Hamid Lone]{\textbf{Dr. Auqib Hamid Lone}}
\institute[GCET Ganderbal]{PCC CS-301: Data Structures\\GCET, Ganderbal}
\date[Week 10]{\small Week 10}

\begin{document}

\frame{\titlepage}

% =============================================================================
% ACT 0 --- MOTIVATION + HANDOFF
% =============================================================================
\section{Motivation and the Handoff}

\begin{frame}{Today's Five Questions}
  \begin{enumerate}
    \item How should records be laid out in a file so that one can be found
      without scanning the rest?
    \item What makes a hash function \emph{good}, and what makes one silently
      bad?
    \item Two keys want the same slot. What are the two families of fixes ---
      and why are their names backwards?
    \item Why does a hash table promise $O(1)$ yet degrade to $O(n)$?
    \item When should a designer refuse hashing and keep order instead?
  \end{enumerate}
  \vspace{0.2cm}
  \begin{exampleblock}{Plan}
    Organize the file. Then replace search with arithmetic: map a key to a
    slot, handle the collisions that arithmetic guarantees, and measure the
    cost with one number --- the load factor $\alpha$.
  \end{exampleblock}
\end{frame}

\begin{frame}{Handoff: What Week 9 Left Open}
  \begin{itemize}
    \item Week 9 sorted the calculator's $2{,}000$-identifier table in
      $\Theta(n\log n)$ and made lookup $\Theta(\log n)$ --- about $11$
      comparisons per identifier.
    \item It closed on one promise: ``Week 10 asks a stronger question: can we
      find a record in \emph{one} probe, without sorting at all? File
      organization and hashing.''
    \item The deficiency it left: $11$ probes is excellent, but sorting is
      paid at startup, and every new identifier must be inserted into a
      sorted array in $\Theta(n)$.
    \item New symbols today: the \key{record} and its \key{key}; the
      \key{hash function} $h(k)$; the \key{load factor} $\alpha = n/m$;
      \key{open hashing} (separate chaining) and \key{closed hashing} (open
      addressing).
  \end{itemize}
  \begin{exampleblock}{This week's job}
    Turn lookup into address arithmetic: $\text{slot} = h(\text{key})$.
    Everything else --- collisions, probing, deletion --- is the price of
    making that arithmetic work on real keys.
  \end{exampleblock}
\end{frame}

\begin{frame}{The Bill for a Lookup}
  \begin{center}
  {\footnotesize
  \begin{tabular}{l r l}
    \toprule
    \textbf{structure for 2{,}000 identifiers} & \textbf{probes} &
      \textbf{insert} \\
    \midrule
    unsorted linked list (Week 6) & $\approx 1{,}000$ & $O(1)$ at head \\
    sorted array + binary search (Weeks 8--9) & $\approx 11$ & $\Theta(n)$ \\
    hash table (this week, $\alpha = 0.5$) & $\approx 1.25$ & $O(1)$ expected \\
    \bottomrule
  \end{tabular}}
  \end{center}
  \begin{itemize}
    \item The sorted array bought a factor of $91$ by \emph{ordering} the
      data. The hash table buys another factor of $\approx 9$ by
      \emph{abandoning} order --- and we must say what that costs.
    \item The hash table does not search for the key at all. It computes where
      the key must be, then checks that one place.
  \end{itemize}
  \muted{Running application (Weeks 3--11): one calculator, rebuilt each
    week. Week 10 replaces its identifier list with a hash table --- the same
    $2{,}000$ names, now found in one or two probes instead of a thousand.}
\end{frame}

\transitionslide{Sorted is one kind of fast. Computed is another.}

% =============================================================================
% ACT 1 --- FILE ORGANIZATION
% =============================================================================
\section{File Organization}

\sectiondivider{Section 1}{Files, Records, and Access Schemes}

\begin{frame}{The Record and the Key}
  \begin{block}{Definition --- record, field, key}
    A \key{record} is a fixed-length collection of fields describing one
    entity. One field is designated the \key{key}: its value identifies the
    record and is the only part a search compares. A \key{file} is a
    collection of records stored on a device.
  \end{block}
  \begin{itemize}
    \item Example: a student record of $64$ bytes
      \texttt{\{roll, name[40], programme, year, marks\}}, keyed by
      \texttt{roll}.
    \item The key is not always unique (names repeat). A \key{primary key} is
      unique; a \key{secondary key} is not, and may index many records.
    \item Every organization in this section answers one question: given a
      key, how many device reads before the record is in memory?
  \end{itemize}
\end{frame}

\begin{frame}{Four Ways to Reach a Record}
  \begin{center}
  {\footnotesize
  \begin{tabular}{l p{4.3cm} p{5.4cm}}
    \toprule
    \textbf{access scheme} & \textbf{how the device is used} &
      \textbf{typical medium} \\
    \midrule
    serial & read every record from the beginning & magnetic tape \\
    sequential & read the next record, in stored order & tape, sorted file \\
    direct (random) & seek straight to a computed block & disk, SSD \\
    indexed & look up a key in an index, then seek & disk, database \\
    \bottomrule
  \end{tabular}}
  \end{center}
  \begin{itemize}
    \item The scheme is a property of the \emph{device and layout}, not of the
      algorithm. A tape cannot do direct access; a disk pays for every seek.
    \item The file organizations that follow are four ways to combine these
      schemes with a record layout. The cost of a lookup is the number of
      \emph{block transfers}, not CPU instructions.
  \end{itemize}
\end{frame}

\begin{frame}{Sequential File}
  \begin{itemize}
    \item Records are stored one after another in the order they arrive
      (\key{pile}) or in key order (\key{sorted sequential}).
    \item \textbf{Unsorted pile.} Search is $\Theta(n)$: scan until found.
      Insert is $\Theta(1)$: append. Delete is $\Theta(n)$ if the gap must be
      closed, or a mark bit.
    \item \textbf{Sorted sequential.} On a tape, search is still a scan; on a
      direct-access device, fixed-length records allow binary search in
      $\Theta(\log n)$ seeks. Insert and delete are $\Theta(n)$ because the
      later records shift.
    \item \textbf{When it wins:} batch jobs that read every record anyway
      (payroll, end-of-term processing), and data too large for any index to
      help.
  \end{itemize}
  \begin{exampleblock}{Worked example}
    $1{,}000$ student records, $64$ bytes each. Finding the last record in an
    unsorted sequential file costs up to $1{,}000$ reads. Sorting the file
    makes a fixed-length binary search possible: $\lceil\log_2 1000\rceil =
    10$ seeks --- but inserting one student now moves the records after it.
  \end{exampleblock}
\end{frame}

\begin{frame}{Direct (Relative) File}
  \begin{itemize}
    \item The record's location is computed from its key, not searched for:
      \[
        \text{offset} = (\text{key} - k_0)\times L,
        \qquad \text{address} = \text{base} + \text{offset},
      \]
      where $k_0$ is the smallest key and $L$ the record length.
    \item Search, insert and delete are all $\Theta(1)$: one seek, one write.
    \item The catch is the key space. Keys must be dense integers. A file for
      employee numbers $\{1, 7, 100{,}000\}$ needs $100{,}000$ slots for three
      records.
  \end{itemize}
  \begin{exampleblock}{Worked example}
    Roll numbers $20230001$--$20231000$, records of $64$ bytes. Roll
    $20230457$ sits at offset $(457-1)\times 64 = 29{,}184$ bytes from the
    base. One seek; no search at all. This is the ideal that hashing will
    imitate for arbitrary keys.
  \end{exampleblock}
\end{frame}

\begin{frame}{Direct File, Drawn}
  \begin{center}
    % Diagram: key -> relative record number -> byte offset -> one seek.
    % Four boxes at y=0, x-centres 0, 2.7, 4.9, 7.4. Arrow labels sit above
    % the arrows, never on them.
    \begin{tikzpicture}[every node/.style={font=\scriptsize}]
      \node[draw, fill=light-bg, minimum width=2.1cm, minimum height=0.9cm,
        align=center] (k) at (0, 0) {roll\\ $20230457$};
      \node[draw, fill=light-bg, minimum width=1.3cm, minimum height=0.9cm,
        align=center] (r) at (3.4, 0) {relative\\ $456$};
      \node[draw, fill=light-bg, minimum width=1.5cm, minimum height=0.9cm,
        align=center] (o) at (5.8, 0) {offset\\ $29{,}184$ B};
      \node[draw, fill=primary-gold!20, minimum width=2.4cm,
        minimum height=0.9cm, align=center] (a) at (8.7, 0) {seek once\\ to base $+$ offset};
      \draw[->, thick, primary-blue] (k) -- (r)
        node[midway, above, text=neutral] {\tiny $-20230001$};
      \draw[->, thick, primary-blue] (r) -- (o)
        node[midway, above, text=neutral] {\tiny $\times 64$};
      \draw[->, thick, primary-blue] (o) -- (a)
        node[midway, above, text=neutral] {\tiny $+$ base};
    \end{tikzpicture}
  \end{center}
  \vspace{-0.1cm}
  \begin{itemize}
    \item The whole organization is one arithmetic pipeline. It works until
      the key space is sparse --- which is the normal case for names, account
      numbers and arbitrary identifiers.
  \end{itemize}
\end{frame}

\begin{frame}{Indexed Sequential File}
  \begin{itemize}
    \item Keep the data file sorted by key, but store it in \emph{blocks}.
      Build a sparse \key{index}: one entry per block, holding that block's
      highest key and its address.
    \item Search: binary-search the index for the block, then scan the block.
      With $B$ blocks of $b$ records, that is $\Theta(\log B + b)$ reads.
    \item Insert: place the record in the correct block if there is room,
      otherwise in an \key{overflow area}; the index still points at the
      block. Periodic \key{reorganization} merges the overflow.
    \item \textbf{When it wins:} point lookups \emph{and} ordered or range
      scans. It is the classic ISAM organization, and the shape a B-tree index
      will later generalize in Week 12's databases.
  \end{itemize}
\end{frame}

\begin{frame}{Hashed File}
  \begin{itemize}
    \item Generalize the direct file to arbitrary keys: compute the bucket
      address instead of deriving it from a dense key.
      \[
        \text{bucket} = h(\text{key}) \bmod m,
        \qquad \text{read that bucket, then search inside it.}
      \]
    \item If two keys hash to the same bucket, the bucket holds more than one
      record --- an \key{overflow chain}. The bucket is the unit of one device
      read.
    \item Search is expected $\Theta(1)$ when the buckets are short. Order is
      gone: there is no ``next record in key order''.
    \item Every remaining section is about making $h$ good and keeping the
      buckets short without destroying the $\Theta(1)$ promise.
  \end{itemize}
  \begin{block}{Definition --- hashed file}
    A file whose record address is $h(\text{key})$; collisions are resolved
    by chaining or by probing. It supports point queries, not range queries.
  \end{block}
\end{frame}

\begin{frame}{The Four Organizations, Side by Side}
  \begin{center}
  {\scriptsize\setlength{\tabcolsep}{4pt}
  \begin{tabular}{p{2.6cm} p{2.1cm} p{2.1cm} p{2.6cm} p{2.5cm}}
    \toprule
    \textbf{organization} & \textbf{search} & \textbf{insert} &
      \textbf{range scan} & \textbf{extra space} \\
    \midrule
    sequential (pile) & $\Theta(n)$ & $\Theta(1)$ & $\Theta(n)$ &
      none \\
    sequential (sorted) & $\Theta(\log n)$ seeks & $\Theta(n)$ &
      $\Theta(n)$ & none \\
    direct / relative & $\Theta(1)$ & $\Theta(1)$ & $\Theta(\text{range})$ &
      key range \\
    indexed sequential & $\Theta(\log B + b)$ & overflow $+$ reorg &
      $\Theta(\log B + \text{out})$ & index $+$ overflow \\
    hashed & $\Theta(1)$ expected & $\Theta(1)$ expected &
      $\Theta(n+m)$ (no order) & table $+$ chain \\
    \bottomrule
  \end{tabular}}
  \end{center}
  \begin{itemize}
    \item Read the table by \emph{which column is fast}. Hashing wins the
      first two columns and loses the third; indexed sequential is the only
      one strong in both.
    \item The choice is a workload question, exactly as it was for sorting in
      Week 9.
  \end{itemize}
\end{frame}

\transitionslide{A hashed file is an array we never search.}

% =============================================================================
% ACT 2 --- HASH FUNCTIONS
% =============================================================================
\section{Hash Functions}

\sectiondivider{Section 2}{Hash Functions and the Load Factor}

\begin{frame}{The Hash-Table ADT}
  \begin{block}{Definition --- hash table (Sedgewick Sec.~3.4, p.~458;
    Wirth Sec.~5.1, p.~177)}
    A hash table stores $n$ keys in $m$ slots. A \key{hash function}
    $h: U \to \{0,\dots,m-1\}$ maps each key to a slot. It supports
    \texttt{insert}, \texttt{search} and \texttt{delete}; the
    \key{load factor} is $\alpha = n/m$.
  \end{block}
  \begin{itemize}
    \item Collisions are unavoidable: the universe $U$ is larger than $m$, so
      two keys must share a slot (pigeonhole). The design question is what to
      do when they do.
    \item The cost model is the \emph{expected number of probes} under
      \key{uniform hashing}: each key is equally likely to hash to any slot,
      independently of the others.
    \item The hash table does not store keys in order. \texttt{search} is
      fast; \texttt{min}, \texttt{max} and ``all keys between $a$ and $b$''
      are not supported.
  \end{itemize}
\end{frame}

\begin{frame}{What Makes a Good Hash Function}
  \begin{itemize}
    \item \textbf{Deterministic.} The same key must always hash to the same
      slot, or a stored record can never be found again.
    \item \textbf{Uniform.} Over the keys actually used, the outputs should be
      spread evenly over the $m$ slots. A function that is uniform on random
      keys but biased on real keys is a bad function.
    \item \textbf{Fast.} Computing $h$ must cost less than the search it
      replaces. A cryptographic hash is uniform but far too slow for a symbol
      table.
    \item \textbf{Uses the whole key.} If the table size is $m = 2^p$, then
      $h(k)=k \bmod m$ sees only the low $p$ bits. Two keys differing only in
      high bits collide every time.
  \end{itemize}
  \begin{exampleblock}{The test that matters}
    Hash the \emph{actual} key distribution, not a textbook's random one.
    Student roll numbers share a year prefix; employee numbers share a
    company prefix; names share letters. A good $h$ must break those patterns.
  \end{exampleblock}
\end{frame}

\begin{frame}{The Division Method}
  \[
    h(k) = k \bmod m.
  \]
  \begin{itemize}
    \item The simplest method, and the fastest. The only design decision is
      $m$ --- and it matters more than the formula suggests.
    \item Choose $m$ \key{prime}, preferably not close to a power of $2$ or
      $10$. Then every digit of $k$ can affect the residue.
    \item Why not $m = 2^p$? Because $k \bmod 2^p$ is exactly the low $p$
      bits: even keys use only half the slots, and sequential keys give
      sequential residues. Why not $m = 10^p$? Decimal keys share trailing
      zeros, so the residue repeats the key's pattern.
  \end{itemize}
  \begin{exampleblock}{Worked example}
    $m = 11$: $h(12)=1$, $h(44)=0$, $h(88)=0$, $h(23)=1$, $h(94)=6$.
    Two collisions already --- exactly what the next two sections handle.
  \end{exampleblock}
\end{frame}

\begin{frame}{The Hash Function Is a Map}
  \begin{center}
    % Diagram: four keys enter h(k)=k mod 11; three slots lie to the right.
    % Key boxes at x=0, y = 1.5, 0.7, -0.1, -0.9; function box centred at
    % (2.3,0.3); slots at x=4.9, y = 1.0, 0.2, -0.6.
    \begin{tikzpicture}[every node/.style={font=\scriptsize}]
      \foreach \y/\v in {1.5/12, 0.7/44, -0.1/88, -0.9/13}
        \node[draw, fill=light-bg, minimum width=1.0cm,
          minimum height=0.5cm] at (0, \y) {\v};
      \node[draw, fill=primary-gold!20, minimum width=2.0cm,
        minimum height=1.6cm, align=center] (h) at (2.3, 0.3)
        {$h(k)$\\ $=k \bmod 11$};
      \foreach \y in {1.5, 0.7, -0.1, -0.9}
        \draw[->, thick, primary-blue] (0.55, \y) -- (1.25, \y);
      \foreach \y/\v in {1.0/0, 0.2/1, -0.6/2}
        \node[draw, circle, fill=light-bg, inner sep=3pt] at (4.9, \y) {\v};
      \draw[->, thick, primary-blue] (3.35, 1.0) -- (4.55, 1.0);
      \draw[->, thick, primary-blue] (3.35, 0.7) -- (4.55, 1.0);
      \draw[->, thick, primary-blue] (3.35, 0.3) -- (4.55, 0.2);
      \draw[->, thick, primary-blue] (3.35, -0.1) -- (4.55, -0.6);
      \draw[dashed, primary-gold, rounded corners]
        (4.55, 0.62) rectangle (5.25, 1.38);
      \node[text=primary-gold, anchor=west] at (5.4, 0.75)
        {\tiny $44$ and $88$ collide};
    \end{tikzpicture}
  \end{center}
  \vspace{-0.1cm}
  \begin{itemize}
    \item Four keys in, three slots out. The collision at slot $0$ is not a
      bug in the arithmetic --- it is the arithmetic working as defined.
  \end{itemize}
\end{frame}

\begin{frame}{The Multiplication Method}
  \[
    h(k) = \left\lfloor m\,(kA \bmod 1)\right\rfloor,
    \qquad 0 < A < 1.
  \]
  \begin{itemize}
    \item Multiply the key by a fixed fraction $A$, keep the fractional part,
      and scale it to $m$ slots. Knuth's choice is
      $A = (\sqrt{5}-1)/2 \approx 0.6180339887$.
    \item Advantage: $m$ need not be prime and is often a power of two, so the
      scaling is a shift. The method uses the \emph{high} bits of the product,
      so it does not ignore the key's leading digits.
    \item Cost: one multiplication and a fractional-part extraction --- still
      constant time, slightly more work than a division.
  \end{itemize}
  \begin{exampleblock}{Worked example}
    $m = 16$, $A = 0.6180339887$. For $k = 12{,}345$,
    $kA = 7{,}629.6296\ldots$, fractional part $0.6296\ldots$, so
    $h(k) = \lfloor 16 \times 0.6296\ldots\rfloor = 10$. The low bits of $k$
    did not decide the slot; the product did.
  \end{exampleblock}
\end{frame}

\begin{frame}{The Classic Four}
  \begin{center}
  {\footnotesize
  \begin{tabular}{l l p{5.4cm}}
    \toprule
    \textbf{method} & \textbf{formula} & \textbf{note} \\
    \midrule
    division & $k \bmod m$ & choose prime $m$; fast; the default \\
    multiplication & $\lfloor m(kA \bmod 1)\rfloor$ &
      any $m$; uses high bits; Knuth's $A$ \\
    mid-square & middle digits of $k^2$ &
      key digits influence all output digits \\
    folding & add or XOR the key's parts &
      good for fixed-width records and IDs \\
    digit analysis & choose the most uniform digits &
      needs a sample of real keys \\
    \bottomrule
  \end{tabular}}
  \end{center}
  \begin{itemize}
    \item \textbf{Mid-square:} $k = 5678$, $k^2 = 32{,}239{,}684$; the middle
      four digits give $2396$.
    \item \textbf{Folding:} $123456789 \to 123 + 456 + 789 = 1368$, then
      $1368 \bmod m$.
    \item \textbf{Digit analysis:} if every roll number begins $2023$, those
      four digits carry no information --- drop them and hash the rest.
  \end{itemize}
\end{frame}

\begin{frame}{Hashing Strings}
  \begin{itemize}
    \item Keys are usually strings, not integers. Treat a string as a
      polynomial in a base $B$ and evaluate it modulo $m$ (Horner's rule):
      \[
        h(s) = \Bigl(\cdots\bigl((s_0 B + s_1)B + s_2\bigr)\cdots\Bigr)
        \bmod m.
      \]
    \item $B = 31$ is the common choice: small, prime, and cheap to multiply
      by a shift and a subtract.
    \item In C, the running value must be reduced modulo $m$ at every step,
      or $31^{n}$ overflows before the string ends. Use \texttt{unsigned}
      arithmetic and reduce inside the loop.
  \end{itemize}
  \begin{exampleblock}{Worked example: \texttt{"cat"} with $m=97$, $B=31$}
    Start $h=0$. $h = 0\cdot31 + 99 = 99$;
    $h = 99\cdot31 + 97 = 3166$;
    $h = 3166\cdot31 + 116 = 98{,}262 \equiv 1 \pmod{97}$.
    So \texttt{"cat"} hashes to slot $1$.
  \end{exampleblock}
\end{frame}

\begin{frame}{Universal Hashing}
  \begin{itemize}
    \item A fixed $h$ can always be defeated: an adversary who knows it can
      choose keys that all collide. For a symbol table fed by untrusted input,
      that is a denial-of-service waiting to happen.
    \item Fix: choose the function \emph{randomly at table creation} from a
      family, and keep the choice secret.
  \end{itemize}
  \begin{block}{Definition --- universal family}
    Let $p$ be a prime larger than the key universe, and choose
    $a \in \{1,\dots,p-1\}$, $b \in \{0,\dots,p-1\}$ uniformly at random.
    Then
    \[
      h_{a,b}(k) = \bigl((ak+b) \bmod p\bigr) \bmod m
    \]
    satisfies $\Pr[h_{a,b}(k_1)=h_{a,b}(k_2)] \le 1/m$ for every distinct
    $k_1,k_2$.
  \end{block}
  \begin{itemize}
    \item No fixed input can force a bad distribution; the expected cost is
      $O(1+\alpha)$ over the random choice. Libraries use this when keys are
      attacker-controlled.
  \end{itemize}
\end{frame}

\begin{frame}{A Bad Hash Function, Traced}
  \begin{itemize}
    \item Suppose the key is a student's year of admission followed by a
      sequence number, and we hash with $m=100$ by taking the \emph{last two
      digits}. Sequential IDs spread perfectly; that is fine.
    \item Now suppose $m = 10$ and every key is even. Then $k \bmod 10 \in
      \{0,2,4,6,8\}$: five slots are never used, and the effective load factor
      is $2\alpha$, not $\alpha$.
    \item Or suppose $m=2^{10}=1024$ and keys differ only in the high-order
      year bits. The remainder keeps only the low $10$ bits, so every student
      in a cohort can collide in the same few slots.
  \end{itemize}
  \begin{exampleblock}{Rule of thumb}
    A hash function is only as good as its behaviour on the keys you will
    actually store. Test it on a sample: count the bucket lengths, and look at
    the longest chain. If one bucket holds a large fraction of $n$, change $m$
    or change $h$ --- do not blame the collision policy.
  \end{exampleblock}
\end{frame}

\transitionslide{Two keys can want the same slot. That is normal, not an error.}

% =============================================================================
% ACT 3 --- OPEN HASHING (SEPARATE CHAINING)
% =============================================================================
\section{Open Hashing}

\sectiondivider{Section 3}{Open Hashing: Separate Chaining}

\begin{frame}{Two Names, One Warning}
  \begin{alertblock}{The names are historically confusing --- define both}
    \begin{itemize}
      \item \key{Open hashing} = \key{separate chaining}: colliding keys are
        stored \emph{outside} the table, in a list attached to the slot.
      \item \key{Closed hashing} = \key{open addressing}: every key is stored
        \emph{inside} the table, and collisions are resolved by probing other
        slots.
    \end{itemize}
  \end{alertblock}
  \begin{itemize}
    \item ``Open'' refers to the table being open to an unbounded number of
      entries per slot. ``Closed'' refers to the address space being closed:
      a key must end up in one of the $m$ slots.
    \item The syllabus uses both pairs. Whenever either name appears, give the
      alternative name in the same sentence --- on an exam, an unqualified
      ``open hashing'' is ambiguous.
  \end{itemize}
\end{frame}

\begin{frame}{Separate Chaining: The Idea}
  \begin{itemize}
    \item The table is an array of $m$ \key{heads}. Slot $j$ points to a
      linked list of every key with $h(k)=j$.
    \item Insert: compute $h(k)$, push the key onto that slot's list ---
      $\Theta(1)$. This is Week 6's linked list, one per bucket.
    \item Search: compute $h(k)$, then walk that one list comparing keys.
      The list is short if $h$ is uniform.
    \item Delete: unlink the node --- $\Theta(1)$ after the search. No
      tombstones, no probing.
    \item $\alpha = n/m$ can exceed $1$: the table has $m$ slots but no limit
      on the total number of nodes.
  \end{itemize}
\end{frame}

\begin{frame}{Chaining, Drawn}
  \begin{center}
    % Diagram: m=11 buckets. Non-empty chains only:
    % 0: 44->88->11, 1: 12->23, 2: 13, 5: 16->5, 6: 94->39, 9: 20.
    % Heads at y=-0.85, deeper nodes at -1.6 and -2.35; n=11, alpha=1.0.
    \begin{tikzpicture}[every node/.style={font=\scriptsize}]
      \foreach \i in {0,...,10}
        \node[draw, fill=light-bg, minimum width=0.7cm,
          minimum height=0.5cm] (b\i) at (\i, 0) {\i};
      \foreach \i/\v in {0/44, 1/12, 2/13, 5/16, 6/94, 9/20}
        \node[draw, fill=primary-gold!20, minimum width=0.62cm,
          minimum height=0.45cm] (h\i) at (\i, -0.85) {\v};
      \foreach \i in {0,1,2,5,6,9}
        \draw[->, thick, primary-blue] (b\i) -- (h\i);
      \node[draw, fill=primary-gold!20, minimum width=0.62cm,
        minimum height=0.45cm] (s0) at (0, -1.6) {88};
      \node[draw, fill=primary-gold!20, minimum width=0.62cm,
        minimum height=0.45cm] (s1) at (1, -1.6) {23};
      \node[draw, fill=primary-gold!20, minimum width=0.62cm,
        minimum height=0.45cm] (s5) at (5, -1.6) {5};
      \node[draw, fill=primary-gold!20, minimum width=0.62cm,
        minimum height=0.45cm] (s6) at (6, -1.6) {39};
      \node[draw, fill=primary-gold!20, minimum width=0.62cm,
        minimum height=0.45cm] (t0) at (0, -2.35) {11};
      \foreach \a/\b in {h0/s0, h1/s1, h5/s5, h6/s6, s0/t0}
        \draw[->, thick, primary-blue] (\a) -- (\b);
      \node[text=neutral, anchor=west] at (10.7, 0)
        {$m=11$};
      \node[text=primary-gold, anchor=west] at (10.7, -0.85)
        {$\alpha=11/11=1$};
    \end{tikzpicture}
  \end{center}
  \vspace{-0.15cm}
  \begin{itemize}
    \item Slot $0$ carries a chain of three; five slots are empty. A search
      for $23$ walks slot $1$'s chain: $12$, then $23$ --- two nodes.
  \end{itemize}
\end{frame}

\begin{frame}{Chaining: Search and Insert Traced}
  \begin{itemize}
    \item Table of the previous slide, $m=11$, $h(k)=k\bmod 11$.
    \item \textbf{Search 23.} $h(23)=1$; read head $1$; compare $12$ (no);
      follow to $23$ (yes). Two comparisons.
    \item \textbf{Search 50 (a miss).} $h(50)=6$; walk slot $6$'s chain
      $94, 39$; reach the end. Two comparisons, and the miss is detected
      without scanning the table.
    \item \textbf{Insert 60.} $h(60)=5$; push at the head of slot $5$; the
      chain becomes $60, 16, 5$. One pointer write --- no shifting, no
      probing.
    \item \textbf{Delete 88.} Search slot $0$ for $88$ (second node), unlink
      it. The chain becomes $44, 11$; no other slot is disturbed.
  \end{itemize}
  \begin{exampleblock}{Why deletion is easy here}
    A chain is a real linked list, so removing a node is Week 6's pointer
    surgery. This is the one operation where chaining is unambiguously
    simpler than open addressing.
  \end{exampleblock}
\end{frame}

\begin{frame}{Chaining: The Bill}
  Under uniform hashing, the expected number of key comparisons in a chain
  is:
  \[
    \text{successful} \approx 1 + \frac{\alpha}{2},
    \qquad
    \text{unsuccessful} \approx \alpha.
  \]
  \begin{center}
  {\footnotesize
  \begin{tabular}{r r r}
    \toprule
    $\alpha = n/m$ & successful & unsuccessful \\
    \midrule
    $0.5$ & $1.25$ & $0.5$ \\
    $1.0$ & $1.5$  & $1.0$ \\
    $2.0$ & $2.0$  & $2.0$ \\
    $4.0$ & $3.0$  & $4.0$ \\
    \bottomrule
  \end{tabular}}
  \end{center}
  \begin{itemize}
    \item The successful cost is half the chain because the target is equally
      likely to be anywhere in it. The unsuccessful cost is the whole chain.
    \item Space is $\Theta(n+m)$: $n$ nodes plus $m$ heads. Chaining degrades
      gracefully --- doubling $n$ doubles the cost, not the square.
  \end{itemize}
\end{frame}

\begin{frame}[fragile]{Chaining in C}
  \scriptsize
\begin{verbatim}
typedef struct Entry {
    const char *key; int value; struct Entry *next;
} Entry;
typedef struct { Entry **head; int m; int n; } Table;

unsigned hash(const char *s, int m) {
    unsigned h = 0;
    for (; *s; s++) h = (h * 31u + (unsigned char)*s) % (unsigned)m;
    return h;
}
int table_search(Table *t, const char *key, int *out) {
    Entry *e = t->head[hash(key, t->m)];
    for (; e; e = e->next)
        if (!strcmp(e->key, key)) { if (out) *out = e->value; return 1; }
    return 0;
}
int table_insert(Table *t, const char *key, int value);  /* push at head */
\end{verbatim}
  \vspace{-0.1cm}
  \footnotesize
  \muted{Insert links one \texttt{Entry} at the head and increments
    \texttt{n}. Resize (rehash into $2m$ slots) when $\alpha > 1$.}
\end{frame}

\transitionslide{When the table must live inside the array itself.}

% =============================================================================
% ACT 4 --- CLOSED HASHING (OPEN ADDRESSING)
% =============================================================================
\section{Closed Hashing}

\sectiondivider{Section 4}{Closed Hashing: Open Addressing}

\begin{frame}{Open Addressing: The Idea}
  \begin{itemize}
    \item Store every key \emph{in the table}. No pointers, no nodes, no
      allocation --- just an array of $m$ slots.
    \item On a collision, follow a \key{probe sequence}:
      $h(k), h(k)+1, h(k)+2, \dots$ (mod $m$) until an empty slot is found.
    \item Search follows the same sequence until the key is found or an
      \emph{empty} slot proves it absent. Delete must leave a marker, because
      clearing a slot would break the proof.
    \item $\alpha = n/m \le 1$ by construction. Performance is controlled by
      keeping $\alpha$ below a threshold ($0.5$--$0.75$) and rehashing when
      it is exceeded.
    \item Cache-friendly: the probe sequence walks contiguous memory, which
      is why open addressing often beats chaining in wall-clock time even at
      the same asymptotic cost.
  \end{itemize}
\end{frame}

\begin{frame}{Linear Probing}
  \[
    h(k,i) = \bigl(h'(k) + i\bigr) \bmod m, \qquad i = 0,1,2,\dots
  \]
  \begin{itemize}
    \item Insert: probe $h'(k)$, $h'(k)+1$, \dots{} until an empty slot.
      Search: the same walk until the key or an empty slot.
    \item Trace with $m=11$, $h'(k)=k\bmod 11$, inserting
      $12,44,13,88,23,94$ in that order:
      \[
        \begin{array}{c|cccccc}
          \text{slot} & 0 & 1 & 2 & 3 & 4 & 6 \\
          \hline
          \text{key} & 44 & 12 & 13 & 88 & 23 & 94
        \end{array}
      \]
      $88$ wants slot $0$ (taken), then $1$ (taken), then $2$ (taken), and
      lands at $3$. $23$ wants $1$, then $2,3$, and lands at $4$.
    \item A search for $23$ now costs four probes --- the collision has been
      paid for by every later search, not just by the insert.
  \end{itemize}
\end{frame}

\begin{frame}{Linear Probing, Drawn}
  \begin{center}
    % Diagram: final 11-slot table after inserting 12,44,13,88,23,94.
    % Slots at x=0..10, y=0. The brace below slots 1..4 marks the four probes
    % made by search(23); the brace is 0.55 cm below the slot boxes.
    \begin{tikzpicture}[every node/.style={font=\scriptsize}]
      \foreach \i/\v in {0/44, 1/12, 2/13, 3/88, 4/23, 5/{}, 6/94}
        \node[draw, fill=light-bg, minimum width=0.8cm,
          minimum height=0.62cm] at (\i, 0) {\v};
      \foreach \i/\v in {7/{}, 8/{}, 9/{}, 10/{}}
        \node[draw, fill=white, minimum width=0.8cm,
          minimum height=0.62cm] at (\i, 0) {\v};
      \foreach \i in {0,...,10}
        \node[text=neutral] at (\i, -0.55) {\tiny \i};
      \draw[decorate, decoration={brace, mirror, amplitude=4pt},
        primary-gold] (0.55, -0.85) -- (4.45, -0.85)
        node[midway, below=7pt, text=primary-gold]
        {\tiny search $23$: $1 \to 2 \to 3 \to 4$};
      \node[text=neutral, anchor=west] at (-0.55, 0.55) {%
        \tiny $h'(23)=1$};
      \node[text=primary-gold, anchor=west] at (6.6, 0.75)
        {\tiny four probes, not one};
    \end{tikzpicture}
  \end{center}
  \vspace{-0.2cm}
  \begin{itemize}
    \item The keys $88$ and $23$ were displaced from their home slots, and
      they push every later key further along. That growing run is
      \key{primary clustering}.
  \end{itemize}
\end{frame}

\begin{frame}{Check Your Prediction: The Next Insert}
  \begin{block}{Predict before you read the answer}
    The table holds $44,12,13,88,23,\cdot,94,\cdot,\cdot,\cdot,\cdot$ at
    slots $0$--$10$. Insert $11$, whose home slot is $h'(11)=0$.
  \end{block}
  \begin{itemize}
    \item It probes $0,1,2,3,4$ --- all occupied --- and lands at slot $5$.
      One insert paid five probes.
    \item A search for $11$ now makes six comparisons ($0$ through $5$).
      The cost of the insert has become the cost of every later lookup of
      that key.
    \item That is the difference between linear probing and a chain: the
      chain's insert touches one node, while the probe's insert commits every
      future search to the same walk.
  \end{itemize}
  \muted{Insert cost is paid once; probe cost is paid on every search. The
    next two sections are about keeping both small.}
\end{frame}

\begin{frame}{Primary Clustering}
  \begin{itemize}
    \item Once a run of occupied slots forms, any key that hashes into the run
      must traverse the whole run and extend it. Runs grow faster than the
      keys that created them.
    \item Knuth's expected probe counts for linear probing:
      \[
        \text{successful} \approx \frac12\left(1+\frac{1}{1-\alpha}\right),
        \qquad
        \text{unsuccessful} \approx \frac12\left(1+\frac{1}{(1-\alpha)^2}\right).
      \]
  \end{itemize}
  \begin{center}
  {\footnotesize
  \begin{tabular}{r r r}
    \toprule
    $\alpha$ & successful & unsuccessful \\
    \midrule
    $0.25$ & $1.17$ & $1.39$ \\
    $0.50$ & $1.50$ & $2.50$ \\
    $0.75$ & $2.50$ & $8.50$ \\
    $0.90$ & $5.50$ & $50.5$ \\
    $0.95$ & $10.5$ & $200.5$ \\
    \bottomrule
  \end{tabular}}
  \end{center}
  \begin{itemize}
    \item The unsuccessful column is the warning: at $\alpha=0.9$, a miss costs
      about $50$ probes. Linear probing is safe only well below $0.75$.
  \end{itemize}
\end{frame}

\begin{frame}{Deletion: The Tombstone}
  \begin{itemize}
    \item Delete $13$ from slot $2$ of the trace. If the slot is simply
      cleared, a later search for $88$ probes $0$ (occupied), $1$ (occupied),
      $2$ (\emph{empty}) --- and concludes that $88$ is absent. It is wrong:
      $88$ sits at slot $3$.
    \item The empty slot is the search's stopping condition, so deletion must
      not create one. Mark the slot \key{DELETED} (a \key{tombstone}).
    \item Search treats \texttt{DELETED} as occupied and continues. Insert may
      reuse the first tombstone it meets. When tombstones dominate, rebuild
      the table (\key{rehash}).
  \end{itemize}
  \begin{alertblock}{Rule}
    \texttt{EMPTY} stops a search; \texttt{DELETED} does not. Every slot has
    one of three states, not two.
  \end{alertblock}
\end{frame}

\begin{frame}{Tombstones, Drawn}
  \begin{center}
    % Diagram: slots 0..4 after deleting 13. Search(88) must pass 0,1,2,3.
    % Slot 2 is DELETED (red dashed); the probing path is a brace below.
    \begin{tikzpicture}[every node/.style={font=\scriptsize}]
      \foreach \i/\v in {0/44, 1/12, 3/88, 4/{}}
        \node[draw, fill=light-bg, minimum width=1.0cm,
          minimum height=0.65cm] at (\i, 0) {\v};
      \node[draw=negative, fill=negative!10, dashed, text=negative,
        minimum width=1.0cm, minimum height=0.65cm, align=center]
        at (2, 0) {\tiny DELETED};
      \foreach \i in {0,...,4}
        \node[text=neutral] at (\i, -0.55) {\tiny \i};
      \draw[decorate, decoration={brace, mirror, amplitude=4pt},
        primary-blue] (-0.45, -0.85) -- (3.45, -0.85)
        node[midway, below=7pt, text=primary-blue]
        {\tiny search $88$: probes $0 \to 1 \to 2 \to 3$};
      \node[text=negative, anchor=west] at (5.4, 0.35)
        {\tiny not empty --- keep going};
      \node[text=neutral, anchor=west] at (5.4, -0.25)
        {\tiny clearing slot $2$ would hide $88$};
    \end{tikzpicture}
  \end{center}
  \vspace{-0.2cm}
  \begin{itemize}
    \item The tombstone is a promise: ``a key may have passed through here.''
      It is the closed-hashing price for having no pointers.
  \end{itemize}
\end{frame}

\begin{frame}{Quadratic Probing}
  \[
    h(k,i) = \bigl(h'(k) + c_1 i + c_2 i^2\bigr) \bmod m.
  \]
  \begin{itemize}
    \item The probe offsets are $0, c_1+c_2, 2c_1+4c_2, \dots$ --- not
      consecutive, so a key does not simply extend a run.
    \item It removes \emph{primary} clustering but creates \key{secondary
      clustering}: keys with the same home slot follow the same sequence.
    \item The guarantee is subtler than linear probing's. With $m$ prime and
      offsets $i^2$, only $\lceil m/2\rceil$ slots are reachable; keep
      $\alpha \le 1/2$. With $m = 2^p$ and offsets $i(i+1)/2$, all slots are
      reachable.
    \item Expected probes are lower than linear at the same $\alpha$, but the
      arithmetic is more expensive and the cache behaviour is worse.
  \end{itemize}
\end{frame}

\begin{frame}{Double Hashing}
  \[
    h(k,i) = \bigl(h_1(k) + i\,h_2(k)\bigr) \bmod m.
  \]
  \begin{itemize}
    \item The step size is a second hash of the key, so two keys with the same
      home slot usually take different paths. This is the closest practical
      approximation to uniform hashing.
    \item Requirement: $h_2(k)$ must be \key{coprime to $m$}, or the sequence
      visits only a fraction of the slots. With prime $m$, choose
      $h_2(k) \in \{1,\dots,m-1\}$; with $m=2^p$, force $h_2(k)$ odd.
    \item Common choice: $h_2(k) = 1 + (k \bmod (m-1))$, which is never $0$
      and is coprime to a prime $m$.
  \end{itemize}
  \begin{exampleblock}{Worked example}
    $m=11$, $h_1(k)=k\bmod 11$, $h_2(k)=1+(k\bmod 10)$. For $k=23$:
    $h_1=1$, $h_2=4$, so the probes are
    $1, 5, 9, 2, 6, 10, 3, 7, 0, 4, 8$ --- every slot exactly once.
  \end{exampleblock}
\end{frame}

\begin{frame}{Probe Sequences Compared}
  \begin{center}
    % Diagram: probe order for key 23 (h1=1) under the three schemes.
    % Rows at y=2.2 (linear), 0.8 (quadratic), -0.6 (double); slots x=0..10.
    % Filled slots carry the probe number 1..4.
    \begin{tikzpicture}[every node/.style={font=\scriptsize}]
      \foreach \r/\lab in {2.2/linear, 0.8/quadratic, -0.6/double}
        \node[text=neutral, anchor=east] at (-0.55, \r) {\lab};
      \foreach \p in {1,2,3,4}
        \node[draw, fill=primary-gold!25, minimum width=0.62cm,
          minimum height=0.45cm] at ({1+\p-1}, 2.2) {\p};
      \foreach \p/\s in {1/1, 2/2, 3/5, 4/10}
        \node[draw, fill=primary-gold!25, minimum width=0.62cm,
          minimum height=0.45cm] at (\s, 0.8) {\p};
      \foreach \p/\s in {1/1, 2/5, 3/9, 4/2}
        \node[draw, fill=primary-gold!25, minimum width=0.62cm,
          minimum height=0.45cm] at (\s, -0.6) {\p};
      \foreach \i in {0,...,10}
        \node[text=neutral] at (\i, -1.3) {\tiny \i};
    \end{tikzpicture}
  \end{center}
  \vspace{-0.1cm}
  \begin{itemize}
    \item Steps: linear $+1$; quadratic $+1,+4,+9$; double $+4$ each time.
      All three find a slot, but only double hashing scatters the probes.
  \end{itemize}
\end{frame}

\begin{frame}[fragile]{Open Addressing in C}
  \scriptsize
\begin{verbatim}
typedef struct {
    const char *key; int value; unsigned char state;  /* 0=E, 1=O, 2=D */
} Slot;

int probe(Table *t, const char *key) {
    unsigned i = hash(key, t->m), step = 1 + hash2(key, t->m - 1);
    for (unsigned k = 0; k < t->m; k++) {
        Slot *s = &t->slot[(i + k*step) % t->m];  /* double hashing */
        if (s->state == EMPTY) return -1;         /* absent */
        if (s->state == OCCUPIED && !strcmp(s->key, key))
            return (int)((i + k*step) % t->m);
    }
    return -1;
}
\end{verbatim}
  \vspace{-0.1cm}
  \footnotesize
  \muted{Search stops only on \texttt{EMPTY}; \texttt{DELETED} keeps the
    probe chain alive. Insert reuses the first \texttt{DELETED} slot it
    passes. Rehash when $\alpha$ or the tombstone count crosses its
    threshold.}
\end{frame}

\transitionslide{Expected $O(1)$ is a promise about the load factor.}

% =============================================================================
% ACT 5 --- CHOOSING AND THE CLOSE
% =============================================================================
\section{Choosing and the Close}

\sectiondivider{Section 5}{Comparison, Choice, and the Close}

\begin{frame}{Chaining vs.\ Open Addressing}
  \begin{center}
  {\footnotesize
  \begin{tabular}{p{3.2cm} p{4.4cm} p{4.4cm}}
    \toprule
    & \textbf{open hashing (chaining)} &
      \textbf{closed hashing (open addressing)} \\
    \midrule
    load factor & any $\alpha$; can exceed $1$ & must stay below $1$ \\
    space & $n$ nodes $+$ $m$ heads & $m$ slots, no pointers \\
    deletion & unlink a node & tombstone $+$ periodic rehash \\
    cache & pointer chasing & contiguous probes \\
    clustering & none --- chains are separate & primary / secondary \\
    ordered scan & sort the chains & none \\
    worst case & $\Theta(n)$ & $\Theta(n)$ \\
    \bottomrule
  \end{tabular}}
  \end{center}
  \begin{itemize}
    \item Chaining is easier to reason about and to delete from; open
      addressing is faster in memory but unforgiving of a high $\alpha$.
  \end{itemize}
\end{frame}

\begin{frame}{Choosing: Read the Workload, Then the Structure}
  \begin{center}
  {\footnotesize
  \begin{tabular}{p{6.4cm} p{5.4cm}}
    \toprule
    \textbf{workload} & \textbf{choose} \\
    \midrule
    point lookups, inserts, deletes, unknown keys & chaining, or double hashing \\
    tight memory, cache-sensitive, $\alpha \le 0.5$ & linear probing \\
    range queries or ordered traversal & indexed sequential or a search tree \\
    adversarial or untrusted keys & universal / randomized hashing \\
    static file, dense integer key & direct (relative) file \\
    batch scan of every record & sequential file \\
    \bottomrule
  \end{tabular}}
  \end{center}
  \begin{itemize}
    \item The same discipline as Week 9: state the workload, then justify the
      structure. ``Hashing is faster'' is not a justification --- it is a
      claim about one column of the table.
    \item The one thing hashing cannot buy is \emph{order}. Week 11's trees
      are the structures that keep both logarithmic search and order.
  \end{itemize}
\end{frame}

\begin{frame}{Three Classic Mistakes}
  \begin{itemize}
    \item \bad{``Open hashing means open addressing.''} It is the reverse.
      Open hashing is separate chaining; closed hashing is open addressing.
      Give both names every time the pair appears.
    \item \bad{``A hash lookup is $O(1)$'' with no qualifier.} It is
      $O(1+\alpha)$ \emph{expected} under uniform hashing, and $\Theta(n)$
      worst case when everything collides. At $\alpha=0.9$, linear probing
      needs about $50$ probes for a miss.
    \item \bad{``To delete, clear the slot.''} An empty slot is the search's
      stopping condition; clearing it hides every key that probed past it.
      Mark \texttt{DELETED}, and rehash when tombstones accumulate.
  \end{itemize}
  \begin{exampleblock}{How to revise}
    For each scheme, write one line: the collision rule, the expected probes
    at $\alpha=0.5$, how deletion works, and whether order survives. Then
    trace an insert and a delete by hand on a table of $11$ slots. The exam
    tests traces and cases, not definitions alone.
  \end{exampleblock}
\end{frame}

\begin{frame}{Week 10 Summary}
  \begin{itemize}
    \item \key{File organization} --- sequential, direct, indexed sequential
      and hashed layouts trade lookup speed against order, space and update
      cost; the access scheme describes the device, the layout decides the
      algorithm.
    \item \key{Hashing} --- $n$ keys in $m$ slots via $h(k)$; $\alpha=n/m$
      sets the expected cost. A good $h$ is deterministic, uniform on the real
      keys, fast, and uses the whole key.
    \item \key{Open hashing} (chaining) --- collisions go in per-slot lists:
      $1+\alpha/2$ probes for a hit, $\alpha$ for a miss; simple deletion;
      $\alpha$ may exceed $1$.
    \item \key{Closed hashing} (open addressing) --- probes inside the table:
      linear, quadratic, double. Deletion needs tombstones; keep $\alpha$
      below $0.5$--$0.75$.
  \end{itemize}
  \begin{exampleblock}{Next week}
    Hashing finds a key in one probe but cannot answer ``the next key'' or
    ``all keys between $a$ and $b$''. Week 11 builds trees --- binary trees,
    BSTs and AVL rotations --- that keep order \emph{and} logarithmic search.
  \end{exampleblock}
  \vspace{-0.2cm}
\end{frame}

\begin{frame}{Before You Go: Predict}
  \begin{block}{Three one-line predictions}
    Answer from the algorithms, then check in the lab.
  \end{block}
  \begin{enumerate}
    \item $m=11$, $h(k)=k\bmod 11$, inserting $12,44,13,88,23,94$ with linear
      probing: where does $88$ land, and how many probes does a search for
      $23$ take?
    \item With linear probing at $\alpha=0.9$, roughly how many probes does an
      unsuccessful search need? What if the table used chaining at the same
      $\alpha$?
    \item You need both point lookups on roll numbers and the query ``list
      all students with roll in $[20231000,20231999]$''. Which file
      organization, and why not a hash table?
  \end{enumerate}
  \muted{Answers: $88$ probes $0,1,2$ and lands at $3$; $23$ probes
    $1,2,3,4$ --- four probes. Linear probing at $0.9$: about $50.5$ probes;
    chaining at $0.9$: about $0.9$ nodes for a miss, so roughly one.
    Indexed sequential (or a tree): hashing destroys order, so a range query
    would have to scan the whole table.}
\end{frame}

\begin{frame}[allowframebreaks]{References}
  \tiny
  \bibliographystyle{plain}
  \bibliography{../../Bibliography_base}
\end{frame}

\end{document}
